\documentclass[11pt,a4paper]{article}
\usepackage[margin=2.5cm]{geometry}
\usepackage{amsmath,amssymb}
\usepackage{booktabs}
\usepackage{enumitem}
\usepackage{xcolor}
\usepackage[colorlinks=true,linkcolor=blue!50!black]{hyperref}

\definecolor{boxbg}{HTML}{EEF3F8}
\newcommand{\idea}[1]{\par\smallskip\noindent\colorbox{boxbg}{\parbox{\dimexpr\linewidth-2\fboxsep}{\textbf{Key idea.} #1}}\par\smallskip}
\newcommand{\where}[1]{\textit{\small (thesis: #1)}}

\title{Reading guide\\[2pt]\large ``The learning problem'' (\texttt{sec:pf:game})}
\author{Companion to \texttt{Chapters/2.Introduction.tex}}
\date{}

\begin{document}
\maketitle

\noindent This guide walks through the subsection \emph{The learning
problem} of the Problem Formulation, one piece at a time. It adds no new
model. Everything here points back to an object the thesis already
defines. Section~\ref{sec:guide-kernel} spends the most time on the
transition kernel $P$, because the thesis only names it and lists where
its parts come from.

\tableofcontents

%----------------------------------------------------------------
\section{What the section is for}

Up to this point the chapter has defined the warehouse piece by piece:
the graph, the agents and how they move, the storage configuration, the
tasks, the controllers and what each one can observe. The learning
problem section does one thing. It shows that all of these pieces
together form a \emph{partially observable stochastic game} (POSG), the
standard framework for several agents that act at the same time, each
seeing the part of the world its controller's window allows. All three
controllers are learned, and Section~\ref{sec:guide-games} shows that they
play three versions of one game that differ only in what each agent
observes.

\idea{Nothing in the POSG is new. It is a repackaging of objects that are
already defined, so that the standard reinforcement-learning machinery
(policies, returns, a training objective) can be applied to them.}

Why bother? Because once the problem is written as a POSG, the training
objective $J(\theta)$ at the end of the section has a precise meaning,
and the methods chapter can use off-the-shelf multi-agent RL ideas
without redefining the environment.

%----------------------------------------------------------------
\section{One timestep in plain words}
\label{sec:guide-loop}

The opening paragraph of the section describes a loop. Here it is with
the symbols attached.

\begin{enumerate}[label=\arabic*.,leftmargin=*]
  \item The warehouse is in some state $s_t$. \where{``realised state''}
  \item Each agent $a_i$ receives its own observation $O_i(s_t)$. Under
        the section-based and decentralised controllers this is only part
        of $s_t$. Under the centralised controller it is all of it.
  \item Every agent picks an action $u_i(t)$ at the same time. Together
        they form the joint action $u_t = (u_1(t),\dots,u_m(t))$.
  \item The warehouse moves to a new state $s_{t+1}$, drawn from the
        transition kernel, $s_{t+1} \sim P(\cdot \mid s_t, u_t)$.
  \item Each agent receives a reward $R_i$ that scores the step.
  \item Repeat until the horizon $T$.
\end{enumerate}

\noindent Two words in the name describe this loop.
\emph{Partially observable} is step~2, the agents do not see all of $s_t$,
except under the centralised controller. \emph{Stochastic game} is
steps~3 and~4, several agents act together and the next state is random.

%----------------------------------------------------------------
\section{The tuple, element by element}
\label{sec:guide-tuple}

The thesis writes the game as
\[
  \mathcal{M} = \bigl( \mathcal{A},\, \mathcal{S},\, \{\mathcal{U}_i\},\,
                       P,\, \{\mathcal{O}_i\},\, \{O_i\},\,
                       \{R_i\},\, \gamma \bigr).
\]
Table~\texttt{tab:posg} in the thesis maps each element to its earlier
definition. The table below adds a plain meaning for each.

\begin{center}
\small
\begin{tabular}{lp{4.3cm}p{6.4cm}}
  \toprule
  Element & Plain meaning & Where it comes from \\
  \midrule
  $\mathcal{A}$ & who acts & the fleet $a_1,\dots,a_m$ \\
  $\mathcal{S}$ & everything that could be true of the warehouse at one instant &
      agent locations $\ell_i(t)$, task sets $\mathcal{Q}_t,\mathcal{B}_t,\mathcal{C}_t$, storage $x_t$, and the conflict-resolution priority order drawn at episode start \\
  $\mathcal{U}_i$ & every action agent $i$ could ever take &
      union of the per-vertex sets $\mathcal{U}(v)$, see Section~\ref{sec:guide-actions} \\
  $P$ & the rules of the warehouse, how one state leads to the next &
      movement, conflicts, task lifecycle, assignment, storage update, order arrivals, see Section~\ref{sec:guide-kernel} \\
  $\mathcal{O}_i$ & every observation agent $i$ could receive & the range of $O_i$, one per architecture \\
  $O_i$ & how agent $i$'s observation is cut out of the state & \texttt{eq:obsfunction}, see Section~\ref{sec:guide-games} \\
  $R_i$ & the score each step earns & deliberately left open, set in the methods chapter \\
  $\gamma$ & how much the future counts & a training parameter \\
  \bottomrule
\end{tabular}
\end{center}

\subsection{Spaces versus realisations}

The section insists on one distinction, and it is worth getting straight
before reading further.

\begin{itemize}[leftmargin=*]
  \item A \emph{space} or \emph{function} is the set of all possibilities
        or the rule. Examples are $\mathcal{S}$, $\mathcal{O}_i$, $O_i$, $P$, $R_i$.
  \item A \emph{realisation} is one concrete value that actually occurs
        in a run. Examples are $s_9$, $o_i(t)$, a reward of $-1$ at $t=9$.
\end{itemize}

\noindent The task example, where $a_2$ is assigned $\tau_1$, makes it concrete. $s_9$ is one realised
state at $t=9$. It records where every agent is (including $a_2$, just
assigned $\tau_1$), that $\tau_1 \in \mathcal{B}_9$, and the storage
configuration $x_9$. $\mathcal{S}$ is the space all such snapshots live
in. The letter clash between $\mathcal{O}_i$ (space) and $O_i$
(function) is exactly why the thesis points to the notation appendix.

%----------------------------------------------------------------
\section{Three games, one difference}
\label{sec:guide-games}

All three controllers are learned policies from one family, trained the
same way. What separates them is only what each agent gets to see. The
thesis says this with one equation, \texttt{eq:obsfunction},
\[
  O_i(s_t) =
  \begin{cases}
    (s_t,\, i) & \text{centralised},\\
    (s^{(q)}_t,\, i) \text{ with } \ell_i(t) \in Y_q & \text{section-based},\\
    o_i(t) & \text{decentralised}.
  \end{cases}
\]

\idea{Everything else in $\mathcal{M}$ is identical across the three.
Same agents, same states, same actions, same $P$, same rewards, same
$\gamma$. So the three learned controllers play three versions of one game
that differ only in $O_i$, and that is exactly the contrast RQ2 asks
about.}

\noindent Read the three cases in plain words.
\begin{itemize}[leftmargin=*]
  \item \textbf{Centralised.} Agent $a_i$'s head sees the whole state plus
        its own index $i$, so it knows which agent it is choosing for. This
        game is fully observed, which is a special case a POSG already
        covers, not a different framework.
  \item \textbf{Section-based.} Agent $a_i$'s head sees the zone it is in,
        plus the halo the boundary protocol passes across. When the agent
        crosses into another zone, its observation switches with it.
  \item \textbf{Decentralised.} Agent $a_i$ sees only its own $o_i(t)$,
        the three pieces of the observation subsection (field of view,
        direction to target, crowding).
\end{itemize}

\noindent The windows are nested. The decentralised observation can be
computed from the section-based one, since the halo has the same depth as
the field of view, and the section-based one from the centralised one.
Messages do not spoil this, because they are not part of any observation.
They belong to the decentralised policy, see
Section~\ref{sec:guide-messages}.
The equation says what a controller \emph{may} condition on. Whether the
realised network uses all of it is a methods question, still open for the
centralised network (the pooling point in \texttt{sec:method:model}).

%----------------------------------------------------------------
\section{The transition kernel \texorpdfstring{$P$}{P}}
\label{sec:guide-kernel}

\subsection{What a transition kernel is}

A transition kernel answers one question.
\begin{quote}
  \emph{If the warehouse is in state $s$ and the agents take joint action
  $u$, how likely is each possible next state $s'$?}
\end{quote}
Formally, $P(s' \mid s, u)$ is a probability distribution over
$\mathcal{S}$ for every pair $(s,u)$, so the probabilities over all
possible $s'$ sum to one. Writing $s_{t+1} \sim P(\cdot \mid s_t, u_t)$
means ``the next state is drawn from that distribution''. The thesis
uses exactly this form in the coupling section,
\[
  s_{t+1} \sim P(\cdot \mid s_t, u_t), \qquad u_t \sim \pi(\cdot \mid s_t).
\]

\idea{$P$ is the environment. The policy chooses $u_t$. Everything that
happens \emph{after} the choice, until the next state, is $P$.}

\subsection{What is inside \texorpdfstring{$P$}{P} in this thesis}

The thesis never writes $P$ as one formula. Table~\texttt{tab:posg}
instead lists what $P$ is made of, \texttt{eq:transition} with conflict
resolution, the shared assignment rule, and \texttt{eq:storageupdate}, and
says it is ``stochastic only via order arrivals''. The thesis also says why
the two shared rules sit in $P$. Neither is an action of any agent, and
both are the same for every controller, so they belong to the environment,
not to a policy. Unpacking that, one step of $P$ does the following, all of
it already defined earlier in the chapter.

\begin{enumerate}[label=(\alph*),leftmargin=*]
  \item \textbf{Move the agents.} Each agent either waits or moves along
        one edge, \texttt{eq:transition}, $\ell_i(t+1) = \ell_i(t)$ or
        $(\ell_i(t),\ell_i(t+1)) \in E$.
  \item \textbf{Resolve conflicts.} If two agents would end on the same
        vertex (\texttt{eq:vertexconflict}) or swap across one edge
        (\texttt{eq:swapconflict}), the conflict-resolution mechanism
        decides who yields. It works through the agents in a priority
        order that is drawn once, at the start of the episode, and then
        kept. That order is part of the state, so the mechanism is
        deterministic given the state and adds no randomness.
  \item \textbf{Update the tasks.} A task moves from $\mathcal{B}_t$ to
        $\mathcal{C}_{t+1}$ when its agent delivers it. The target
        $q_i(t)$ switches from pickup to delivery at pickup.
  \item \textbf{Update storage, but only on the slow clock.} At epochs
        $t \in \{\Delta, 2\Delta,\dots\}$ the storage rule $F$ produces
        a new $x_t$ (\texttt{eq:storageupdate}). Between epochs $x_t$
        stays fixed.
  \item \textbf{Add new orders.} New orders are drawn from
        $P_{\mathrm{ord}}$ and turned into tasks in $\mathcal{Q}_{t+1}$.
        \emph{This is the only random step.}
  \item \textbf{Assign tasks.} Free agents are matched to waiting tasks by
        the one assignment rule every controller shares. No policy chooses
        this, so it is part of the environment. Ties are broken from the
        state alone. Tasks released at the same time go in task order, and
        equally close agents are ranked by the same priority order the
        conflict-resolution step uses.
\end{enumerate}

\noindent The methods chapter runs orders and assignment at the
\emph{start} of a timestep, before routing, and resolution and storage
after. Read as a POSG, nothing changes. The end of one timestep and the
start of the next together make up one application of $P$.

\subsection{Why ``stochastic only via order arrivals'' matters}

Steps (a) to (d) and (f) are deterministic. The modelling assumption
``deterministic execution'' says an accepted action always succeeds, so a
robot told to move does move. The only thing the agents cannot predict
is which orders arrive next. One helpful way to read $P$ is therefore in
two parts,
\[
  s_{t+1} = f\bigl(s_t,\, u_t,\, z_{t+1}\bigr), \qquad
  z_{t+1} \sim P_{\mathrm{ord}},
\]
where $f$ collects steps (a) to (d), turning the new orders $z_{t+1}$
into tasks, and step (f). This split is a reading aid, not notation the
thesis uses. It shows that if you fixed the order stream, the whole run
would be completely determined by the agents' actions.

\paragraph{Small example.} Take the running example at $t=9$. Agent
$a_2$ has just been assigned $\tau_1$ and heads for pickup vertex $A$.
Suppose $a_2$ chooses to move one vertex towards $A$ and no other agent
wants that vertex. Then steps (a) to (c) fix $a_2$'s next location
exactly. Unless $t=9$ is a storage epoch, $x_{10}=x_9$. The only
uncertainty in $s_{10}$ is which new tasks appear in $\mathcal{Q}_{10}$.
If no orders could arrive, $P(\cdot \mid s_9, u_9)$ would put
probability one on a single next state.

\subsection{Where the storage layer sits}

It can seem odd that the storage rule $F$ is inside the environment,
since storage is one of the two layers the thesis studies. The reading
is that, from the routing policy's point of view, $F$ is a fixed rule
that runs on its own clock. The routing agents do not choose storage.
They experience its effect as part of how the world evolves, through
pickup vertices moving. That is the coupling the next section of the
thesis (\emph{How the two layers close the loop}) builds on.

%----------------------------------------------------------------
\section{Actions and masking}
\label{sec:guide-actions}

A POSG normally has a fixed action set per agent. In a graph, the legal
actions depend on where the agent stands. At vertex $v$ only
$\mathcal{U}(v) = \{\mathrm{wait}\} \cup \{\mathrm{move}(w) : w \in N^+(v)\}$
is legal. The thesis resolves this in the usual way.

\begin{itemize}[leftmargin=*]
  \item Give every agent the big set $\mathcal{U}_i = \bigcup_v \mathcal{U}(v)$,
        all moves to all vertices, plus wait.
  \item At each state, switch off (mask) everything outside
        $\mathcal{U}(\ell_i(t))$.
\end{itemize}

\noindent Example. At a vertex with one outgoing edge to $w$, the agent
really has two choices, wait or $\mathrm{move}(w)$. All other entries
of $\mathcal{U}_i$ are masked.

\noindent The masking happens in the policy, not in $P$. The masked
actions get probability zero because they are removed before the policy
samples (\texttt{eq:mask} in the methods chapter). $P$ only ever assigns
probabilities to next states.

%----------------------------------------------------------------
\section{The objective \texorpdfstring{$J(\theta)$}{J(theta)}}

\[
  \max_{\theta} J(\theta)
    = \max_{\theta} \mathbb{E}_{\pi_\theta}
      \Bigl[ \sum_{t=0}^{T-1} \gamma^{t}\, \frac{1}{m} \sum_{i=1}^{m} R_i(t) \Bigr]
\]
Read it from the inside out.
\begin{itemize}[leftmargin=*]
  \item $\frac{1}{m}\sum_i R_i(t)$ is the fleet's average reward at step
        $t$. Averaging, not summing per agent, is what makes this a
        \emph{team} objective. The agents are cooperative and share one
        set of weights.
  \item $\gamma^t$ with $\gamma \in [0,1)$ makes later reward count less.
        With $\gamma=0.99$, reward 100 steps ahead counts about
        $0.99^{100}\approx 0.37$ as much as reward now.
  \item $\mathbb{E}_{\pi_\theta}$ averages over everything random in a
        run, the policy's own sampling of actions and the order arrivals
        inside $P$.
  \item $\max_\theta$ says training searches for the weights that make
        this expected team return as large as possible.
\end{itemize}

\noindent $R_i(t)$ is shorthand for $R_i(s_t, u_t, s_{t+1})$, the reward
for agent $i$ on the step from $s_t$ to $s_{t+1}$ under joint action $u_t$.

\noindent The same objective covers all three learned controllers. Only
what $\theta$ parameterises changes, the shared per-agent network in the
decentralised case, or the global or zone-level network otherwise. A
classical reference planner is reported in the results as a baseline, but
it is not part of this problem and is not trained.

\paragraph{One objective, different credit.} If all three controllers
maximise the same team average, why would they differ in how much
congestion they cause? Because the objective is the \emph{target}, and
how each update credits reward to an agent's own action is a separate
training choice. An update can use the mean reward over the agents a
controller moves, or each agent's own $R_i$. With its own $R_i$, an agent
is not charged for the delay it causes others, so it has no reason to
avoid it. That choice is fixed in the methods chapter (it is still marked
open there), and it is what makes the predicted ordering, centralised
least exposed to congestion and decentralised most, consistent with one
shared objective.

\paragraph{The critic sees what the policy sees.} PPO trains a critic
next to each policy, a value estimate used only while learning. In this
thesis each critic sees exactly the window of its own policy, so a
decentralised learner also \emph{learns} from local information only.
That is on purpose. It means the architectures differ in training as well
as in acting, and a decentralised learner is never taught the cost of
congestion it causes out of its sight, which is the effect the comparison
measures.

\paragraph{Why $R_i$ is left open.} The reward is a tool for getting a
good policy, not the thing the thesis finally measures. The final
measures are the system-level criteria (throughput, service time and so
on). The reward can therefore be redesigned without changing the
problem, which is why it is fixed in the methods chapter instead.

%----------------------------------------------------------------
\section{Messages belong to the policy}
\label{sec:guide-messages}

A decentralised agent can talk to the agents within its field of view.
What it sends is not chosen by hand. The agent's own network computes a
message from its own observation, and training decides what is worth
sending. This is how CRAMP and PICO do it, and the thesis follows them.

\idea{A learned message depends on the network's weights, not only on the
state. So it cannot be part of the observation, which is a fixed rule
applied to the state. It is part of the policy instead, a communication
step inside the decentralised controller.}

\noindent Three consequences follow.
\begin{itemize}[leftmargin=*]
  \item $o_i(t)$ has three pieces, field of view, direction to target and
        crowding. Messages are not a fourth piece.
  \item The observation functions stay nested exactly. Nothing a
        decentralised agent observes is hidden from the other two
        controllers.
  \item Communication has its own on/off switch, the conditional RQ5. With
        it off, the agent acts from $o_i(t)$ alone.
\end{itemize}

\noindent Messages travel only one hop, from agents within the field of
view. More hops would pass on what agents further away send, which would
let a decentralised agent learn about agents it cannot see and widen its
window.

%----------------------------------------------------------------
\section{Where the congestion feature fits}

$\delta_i(t)$ can be read from every controller's window, and for the
decentralised controller it is one of the three pieces of $o_i(t)$. So the
policy can see it, and a reward designer may use it. Whether the policy is allowed to
react to it is the experimental toggle ``congestion sensitivity:
disabled/enabled''. That toggle is a choice of the experiment, not a
property of any architecture.

%----------------------------------------------------------------
\section{Points worth checking in the thesis text}

An earlier reading of the section for this guide turned up four places
where the wording was looser than the rest of the formulation. Three are
now fixed in the thesis (tag \texttt{[A12]}). The same pass found two
more, and all of them are now fixed.

\paragraph{Fixed.}
\begin{enumerate}[leftmargin=*]
  \item \textbf{Masked actions and $P$.} The text said the masked actions
        ``have zero probability under $P$'', but $P$ assigns probabilities
        to next states, not to actions. It now says they get zero
        probability under the policy, since they are masked out before
        sampling.
  \item \textbf{The signature of $R_i$.} The thesis now states once that
        $R_i(t)$ is shorthand for $R_i(s_t,u_t,s_{t+1})$.
  \item \textbf{$O_i$ and messages.} The worry was that messages come from
        other agents' policies, so $O_i$ might not be a function of $s_t$
        alone. It was first settled by fixing the message to two numbers
        read from the state. It is now settled the other way (tag
        \texttt{[A14]}). Messages are learned, so they are not part of the
        observation at all but part of the decentralised policy, and
        $O_i$ is a function of the state by construction. See
        Section~\ref{sec:guide-messages}.
  \item \textbf{Ties in the assignment rule.} Once assignment is part of
        $P$, $P$ needs it to be a function, and ties are common when every
        edge costs the same. Tasks released at the same time now go in task
        order, and equally close agents are ranked by the
        conflict-resolution priority order. That order is already in the
        state and is redrawn each episode, so no agent is favoured. Using
        communication to settle ties was ruled out, because assignment
        must be the same for every controller and only the decentralised
        one communicates.
  \item \textbf{Which critic each controller uses} (tag \texttt{[A13]}).
        The thesis used to define centralised training with decentralised
        execution without saying whether it used it. It now says each
        controller's critic sees exactly the window its own policy sees,
        and gives the reason. The window should limit how a controller
        learns as well as how it acts, and a local critic keeps the
        congestion contrast measurable.
  \item \textbf{Two running examples.} The section said ``continuing the
        running example'' right after the observation subsection, which has
        its own running example with agent $a_1$, so the phrase pointed at
        the wrong one. It now names the task example of the task-route
        figure, where $a_2$ is assigned $\tau_1$ at $t=9$.
\end{enumerate}

%----------------------------------------------------------------
\section{Summary in five lines}

\begin{enumerate}[leftmargin=*]
  \item The warehouse problem is a POSG, many agents, one shared environment, and three observation functions, one per controller.
  \item Every element of the POSG was already defined earlier in the chapter, and only $O_i$ differs between the three learned controllers.
  \item $P$ is the environment's rulebook, movement, conflict resolution, task updates, assignment, slow storage updates and order arrivals.
  \item Only the order arrivals are random, everything else in $P$ is deterministic.
  \item Training maximises the discounted team average reward, and how each update credits that reward is a separate choice made in the methods chapter.
\end{enumerate}

\end{document}
